\documentclass[10pt]{article}

% Page geometry
\usepackage[margin=1in]{geometry}

% Math packages
\usepackage{amsmath}
\usepackage{amssymb}
\usepackage{amsthm}

% Lists
\usepackage{enumitem}

% Hyperlinks
\usepackage{hyperref}
\hypersetup{
  colorlinks=true,
  linkcolor=blue,
  citecolor=blue,
  urlcolor=blue
}

% Theorem environments
\newtheorem{theorem}{Theorem}[section]
\newtheorem{definition}[theorem]{Definition}
\theoremstyle{remark}
\newtheorem{remark}[theorem]{Remark}

\title{Spectral Graph Theory of Migration Strategies}

\author{
  Robin Langer \\
  \texttt{langer.robin@gmail.com}
  \and
  Claudius Turing \\
  \texttt{11o1111o11oo1o1o@gmail.com}
  \and
  Lyra Vega \\
  \texttt{lyraclaud20@gmail.com}
}

\date{}

\begin{document}

\maketitle

\section{Migration Graph Framework}

\begin{definition}[Migration Graph]
Let $G = (V, E)$ be an undirected graph on $n$ nodes (islands). An edge
$\{i,j\} \in E$ indicates that individuals may migrate directly between
islands $i$ and $j$. The \emph{migration matrix} $M \in [0,1]^{n \times n}$
is a row-stochastic matrix where $M_{ij} > 0$ only if $\{i,j\} \in E$.
\end{definition}

\begin{definition}[Algebraic Connectivity]
For a graph $G$ with Laplacian $L = D - A$ (where $D =
\mathrm{diag}(\deg_1, \ldots, \deg_n)$ and $A$ is the adjacency matrix),
the \emph{algebraic connectivity} $\lambda_2(G)$ is the second-smallest
eigenvalue of $L$.
\end{definition}

\begin{theorem}[Spectral Ordering of Migration Strategies]
\label{thm:spectral}
Let $G_{\mathrm{strict}} = \bar{K}_n$, $G_{\mathrm{ring}} = C_n$,
$G_{\mathrm{star}} = K_{1,n-1}$, $G_{\mathrm{FC}} = K_n$ denote the
migration graphs for the strict, ring, star, and fully-connected
strategies, respectively. For $n \geq 7$:
\[
  0 = \lambda_2(\bar{K}_n)
    \;<\; \lambda_2(C_n)
    \;<\; \lambda_2(K_{1,n-1}) = 1
    \;<\; \lambda_2(K_n) = n.
\]
Under a random-walk model of migration, this spectral ordering predicts
the inverse diversity stratification ordering:
\[
  \mathrm{div}(\bar{K}_n)
  \;>\; \mathrm{div}(C_n)
  \;>\; \mathrm{div}(K_{1,n-1})
  \;>\; \mathrm{div}(K_n).
\]
Our experiments confirm this ordering holds empirically across all domains.
\end{theorem}

\begin{proof}
\textbf{Step 1: Spectral gap values.}
\begin{itemize}
  \item $\bar{K}_n$ (no edges): disconnected, so $\lambda_2 = 0$ by
    definition.
  \item $C_n$ (ring on $n$ nodes): the eigenvalues of $L(C_n)$ are
    $2(1 - \cos(2\pi k/n))$ for $k = 0, 1, \ldots, n-1$. The smallest
    nonzero eigenvalue is
    $\lambda_2(C_n) = 2(1 - \cos(2\pi/n)) \approx 4\pi^2/n^2$
    for large $n$.
  \item $K_{1,n-1}$ (star with hub of degree $n-1$ and $n-1$ leaves of
    degree 1): the eigenvalues of $L(K_{1,n-1})$ are $0$, $1$ (with
    multiplicity $n-2$), and $n$. Hence $\lambda_2(K_{1,n-1}) = 1$.
    (Proof: $\mathrm{tr}(L) = 2(n-1)$ and $0 + (n-2)\cdot 1 + \lambda_{\max} = 2(n-1)$
    gives $\lambda_{\max} = n$.)
  \item $K_n$ (complete graph, $(n-1)$-regular): the eigenvalues of
    $L(K_n)$ are $0$ (once) and $n$ (with multiplicity $n-1$). Hence
    $\lambda_2(K_n) = n$.
\end{itemize}

For $n \geq 7$: $2(1 - \cos(2\pi/n)) < 1$ iff $\cos(2\pi/n) > \tfrac{1}{2}$
iff $n > 6$. The strict ordering $0 < \lambda_2(C_n) < 1 < n$ therefore
holds for all $n \geq 7$.

\textbf{Step 2: Mixing time.}
Since both $C_n$ (for even $n$) and $K_{1,n-1}$ are bipartite, the simple
random walk $P = D^{-1}A$ has period~2 and does not converge. We therefore
work with the lazy walk $P' = (I + P)/2$, which is aperiodic. (In the GA
setting, migration matrices have positive diagonal---individuals that do not
migrate---so aperiodicity holds naturally.)

For a connected, aperiodic chain the mixing time satisfies
\[
  t_{\mathrm{mix}} = \Theta\!\left(\frac{\log(n/\varepsilon)}{\gamma}\right),
\]
where $\gamma$ is the spectral gap of $P'$
\cite{LevinPeresWilmer2009}. For $d$-regular graphs,
$\gamma = \lambda_2(L)/(2d)$ (the factor $1/2$ from the lazy step); for
non-regular graphs, the spectral gap of the lazy walk satisfies
$\gamma \geq \lambda_2(L)/(2\Delta)$, where $\Delta$ is the maximum degree
\cite{LevinPeresWilmer2009}. In all cases, $\gamma$ is monotone increasing
in $\lambda_2(L)$, so slower mixing corresponds to smaller $\lambda_2$.

\textbf{Step 3: Mixing time $\Leftrightarrow$ diversity.}
Under neutral drift with random-walk migration, the total variation of the
population distribution from stationarity satisfies
\[
  \max_i \bigl\| P^t(i, \cdot) - \pi \bigr\|_{\mathrm{TV}}
  \;\leq\; \tfrac{1}{2}(1 - \gamma)^t.
\]
Population differentiation (e.g., Wright's $F_{ST}$) decays at the same
rate: smaller $\gamma$ $\Rightarrow$ slower equilibration across islands
$\Rightarrow$ higher sustained genetic diversity. Thus, under this model,
diversity is monotone decreasing in $\gamma$ (and hence in $\lambda_2(L)$
for regular graphs).

\textbf{Step 4: Combining.}
The spectral ordering from Step~1 and the monotonicity from Steps~2--3
give, for $n \geq 7$:
\[
  \mathrm{div}(\bar{K}_n)
  \;>\; \mathrm{div}(C_n)
  \;>\; \mathrm{div}(K_{1,n-1})
  \;>\; \mathrm{div}(K_n).
\]
\end{proof}

\begin{remark}[Optimality of the ring among sparse topologies]
Among connected graphs on $n$ nodes with at most $n$ edges (cycles and
unicyclic graphs), $C_n$ achieves the minimum algebraic connectivity
$\lambda_2 \approx 4\pi^2/n^2$. By comparison, any tree on $n$ nodes
satisfies $\lambda_2 \leq 2(1 - \cos(\pi/n)) \approx 2\pi^2/n^2$
(achieved by the path $P_n$), but trees have only $n-1$ edges and are not
2-regular. The ring is therefore the diversity-optimal migration topology
among all connected topologies with bounded degree~2: it mixes the slowest,
preserving the most inter-island differentiation.
\end{remark}

\begin{thebibliography}{1}
\bibitem{LevinPeresWilmer2009}
D.~A. Levin, Y.~Peres, and E.~L. Wilmer,
\textit{Markov Chains and Mixing Times},
American Mathematical Society, 2009.
\end{thebibliography}

\end{document}